% !TEX program = xelatex
\documentclass[12pt]{article}
\usepackage[paperwidth=612bp,paperheight=792bp,margin=0bp,noheadfoot]{geometry}
\usepackage{fontspec}
\usepackage{xeCJK}
\usepackage{tikz}
\pagestyle{empty}
\setlength{\parindent}{0pt}
\setlength{\parskip}{0pt}
\IfFontExistsTF{TeX Gyre Termes}{\setmainfont{TeX Gyre Termes}}{\setmainfont{Times New Roman}}
\IfFontExistsTF{Noto Serif CJK SC}{\setCJKmainfont{Noto Serif CJK SC}}{\setCJKmainfont{Songti SC}}
\newcommand{\QuestionTitle}{人机实验1 H-C5b：项羽本纪}
\newcommand{\DrawQuestionTitle}{%
\node[anchor=base west,inner sep=0pt,outer sep=0pt] at ([xshift=72bp,yshift=-34bp]current page.north west) {{\fontsize{14bp}{14bp}\selectfont\bfseries \QuestionTitle}};%
\node[anchor=base west,inner sep=0pt,outer sep=0pt] at ([xshift=72bp,yshift=-62bp]current page.north west) {{\fontsize{12bp}{12bp}\selectfont 用时：\underline{\hspace{80bp}}}};%
}
\begin{document}
% Auto-generated from 项羽_chapter_with_numbers.tex
\thispagestyle{empty}
\begin{tikzpicture}[remember picture,overlay]
\DrawQuestionTitle
\node[anchor=base west,inner sep=0pt,outer sep=0pt] at ([xshift=72bp,yshift=-84bp]current page.north west) {{\fontsize{12bp}{12bp}\selectfont 【7.8.9】 沛公、项羽离开外黄攻打陈留，陈留坚兵固守，不能攻下。}};
\node[anchor=base west,inner sep=0pt,outer sep=0pt] at ([xshift=72bp,yshift=-100bp]current page.north west) {{\fontsize{12bp}{12bp}\selectfont 【7.8.10】 沛公、项羽互相商量说：“如今项梁的军队垮了，士卒恐惧。”}};
\node[anchor=base west,inner sep=0pt,outer sep=0pt] at ([xshift=72bp,yshift=-116bp]current page.north west) {{\fontsize{12bp}{12bp}\selectfont 【7.8.11】 于是就领兵同吕臣的军队一起向东进发。}};
\node[anchor=base west,inner sep=0pt,outer sep=0pt] at ([xshift=72bp,yshift=-132bp]current page.north west) {{\fontsize{12bp}{12bp}\selectfont 【7.8.12】 吕臣驻扎在彭城东面，项羽驻扎在彭城西面，沛公驻扎在砀。}};
\node[anchor=base west,inner sep=0pt,outer sep=0pt] at ([xshift=72bp,yshift=-148bp]current page.north west) {{\fontsize{12bp}{12bp}\selectfont 【7.9.1】 章邯已经打垮了项梁的军队，以为楚地的敌人不用担心了，就渡过黄河攻打赵}};
\node[anchor=base west,inner sep=0pt,outer sep=0pt] at ([xshift=72bp,yshift=-164bp]current page.north west) {{\fontsize{12bp}{12bp}\selectfont 地，大破赵军。}};
\node[anchor=base west,inner sep=0pt,outer sep=0pt] at ([xshift=72bp,yshift=-180bp]current page.north west) {{\fontsize{12bp}{12bp}\selectfont 【7.9.2】 这个时候，赵歇为赵王，陈余为将，张耳为相，都跑进了巨鹿城。}};
\node[anchor=base west,inner sep=0pt,outer sep=0pt] at ([xshift=72bp,yshift=-196bp]current page.north west) {{\fontsize{12bp}{12bp}\selectfont 【7.9.3】 章邯命令王离、涉间围攻巨鹿，章邯驻扎在巨鹿南面，修筑甬道输送粮食。}};
\node[anchor=base west,inner sep=0pt,outer sep=0pt] at ([xshift=72bp,yshift=-212bp]current page.north west) {{\fontsize{12bp}{12bp}\selectfont 【7.9.4】 陈余作为将领，统率士卒数万人驻扎在巨鹿的北面，这就是所说的河北之军。}};
\node[anchor=base west,inner sep=0pt,outer sep=0pt] at ([xshift=72bp,yshift=-228bp]current page.north west) {{\fontsize{12bp}{12bp}\selectfont 【7.10.1】 楚军在定陶打了败仗，楚怀王很恐惧，从盱台前往彭城，合并了项羽、吕臣的}};
\node[anchor=base west,inner sep=0pt,outer sep=0pt] at ([xshift=72bp,yshift=-244bp]current page.north west) {{\fontsize{12bp}{12bp}\selectfont 军队亲自统率。}};
\node[anchor=base west,inner sep=0pt,outer sep=0pt] at ([xshift=72bp,yshift=-260bp]current page.north west) {{\fontsize{12bp}{12bp}\selectfont 【7.10.2】 以吕臣为司徒，用他的父亲叶青为令尹。}};
\node[anchor=base west,inner sep=0pt,outer sep=0pt] at ([xshift=72bp,yshift=-276bp]current page.north west) {{\fontsize{12bp}{12bp}\selectfont 【7.10.3】 以沛公为砀郡长，封为武安侯，统率砀郡的军队。}};
\node[anchor=base west,inner sep=0pt,outer sep=0pt] at ([xshift=72bp,yshift=-292bp]current page.north west) {{\fontsize{12bp}{12bp}\selectfont 【7.11.1】 以前宋义所遇到的齐国使者高陵君显还在楚国的军队里，他见到楚怀王说：“}};
\node[anchor=base west,inner sep=0pt,outer sep=0pt] at ([xshift=72bp,yshift=-308bp]current page.north west) {{\fontsize{12bp}{12bp}\selectfont 宋义断定武信君的军队一定失败，过了几天，他的军队果然失败了。军队没有开战而先看}};
\node[anchor=base west,inner sep=0pt,outer sep=0pt] at ([xshift=72bp,yshift=-324bp]current page.north west) {{\fontsize{12bp}{12bp}\selectfont 到了失败的征兆，这可说是懂得军事了。”}};
\node[anchor=base west,inner sep=0pt,outer sep=0pt] at ([xshift=72bp,yshift=-340bp]current page.north west) {{\fontsize{12bp}{12bp}\selectfont 【7.11.2】 楚怀王召见宋义，和他商量事情，大为高兴，因此委任为上将军，项羽为鲁公}};
\node[anchor=base west,inner sep=0pt,outer sep=0pt] at ([xshift=72bp,yshift=-356bp]current page.north west) {{\fontsize{12bp}{12bp}\selectfont ，担任次将，范增为末将，去援救赵国。}};
\node[anchor=base west,inner sep=0pt,outer sep=0pt] at ([xshift=72bp,yshift=-372bp]current page.north west) {{\fontsize{12bp}{12bp}\selectfont 【7.16.3】 等到秦军投降了诸侯军，诸侯军的官兵乘战争胜利的机会，像对待奴隶和俘虏}};
\node[anchor=base west,inner sep=0pt,outer sep=0pt] at ([xshift=72bp,yshift=-388bp]current page.north west) {{\fontsize{12bp}{12bp}\selectfont 一样地驱使他们，随便折磨侮辱秦军官兵。}};
\node[anchor=base west,inner sep=0pt,outer sep=0pt] at ([xshift=72bp,yshift=-404bp]current page.north west) {{\fontsize{12bp}{12bp}\selectfont 【7.16.4】 秦军官兵多在私下议论说：“章将军等欺骗我们投降诸侯军。如今能够入关破}};
\node[anchor=base west,inner sep=0pt,outer sep=0pt] at ([xshift=72bp,yshift=-420bp]current page.north west) {{\fontsize{12bp}{12bp}\selectfont 秦，（当然）很好；如果不能，诸侯军俘虏我们东去；秦势必把我们的父母妻子全部处死}};
\node[anchor=base west,inner sep=0pt,outer sep=0pt] at ([xshift=72bp,yshift=-436bp]current page.north west) {{\fontsize{12bp}{12bp}\selectfont 。”}};
\node[anchor=base west,inner sep=0pt,outer sep=0pt] at ([xshift=72bp,yshift=-452bp]current page.north west) {{\fontsize{12bp}{12bp}\selectfont 【7.16.5】 诸侯军的将领们暗中听到了他们的打算，报告了项羽。}};
\node[anchor=base west,inner sep=0pt,outer sep=0pt] at ([xshift=72bp,yshift=-468bp]current page.north west) {{\fontsize{12bp}{12bp}\selectfont 【7.16.6】 项羽就找来黥布、蒲将军商量说：“秦军官兵还很多，他们心里不服，到了关}};
\node[anchor=base west,inner sep=0pt,outer sep=0pt] at ([xshift=72bp,yshift=-484bp]current page.north west) {{\fontsize{12bp}{12bp}\selectfont 中不听从命令，事情必然岌岌可危，不如杀掉他们，而只与章邯、长史司马欣、都尉董翳}};
\node[anchor=base west,inner sep=0pt,outer sep=0pt] at ([xshift=72bp,yshift=-500bp]current page.north west) {{\fontsize{12bp}{12bp}\selectfont 一起入秦。”}};
\node[anchor=base west,inner sep=0pt,outer sep=0pt] at ([xshift=72bp,yshift=-516bp]current page.north west) {{\fontsize{12bp}{12bp}\selectfont 【7.16.7】 于是楚军夜间把秦军士卒二十多万人处死掩埋在新安城南。}};
\node[anchor=base west,inner sep=0pt,outer sep=0pt] at ([xshift=72bp,yshift=-532bp]current page.north west) {{\fontsize{12bp}{12bp}\selectfont 【7.17.1】 项羽将要攻取秦关中地带。}};
\node[anchor=base west,inner sep=0pt,outer sep=0pt] at ([xshift=72bp,yshift=-548bp]current page.north west) {{\fontsize{12bp}{12bp}\selectfont 【7.17.2】 函谷关有兵把守，不能进去。}};
\node[anchor=base west,inner sep=0pt,outer sep=0pt] at ([xshift=72bp,yshift=-564bp]current page.north west) {{\fontsize{12bp}{12bp}\selectfont 【7.17.3】 又听说沛公已经攻破咸阳，项羽大怒，派当阳君等扣关。}};
\node[anchor=base west,inner sep=0pt,outer sep=0pt] at ([xshift=72bp,yshift=-580bp]current page.north west) {{\fontsize{12bp}{12bp}\selectfont 【7.17.4】 项羽便进入了函谷关，到达戏水西岸。}};
\node[anchor=base west,inner sep=0pt,outer sep=0pt] at ([xshift=72bp,yshift=-596bp]current page.north west) {{\fontsize{12bp}{12bp}\selectfont 【7.17.5】 沛公驻军霸上，没有能够和项羽相见。}};
\node[anchor=base west,inner sep=0pt,outer sep=0pt] at ([xshift=72bp,yshift=-612bp]current page.north west) {{\fontsize{12bp}{12bp}\selectfont 【7.17.6】 沛公左司马曹无伤派人对项羽说：“沛公想称王关中，使子婴为相，占有了全}};
\node[anchor=base west,inner sep=0pt,outer sep=0pt] at ([xshift=72bp,yshift=-628bp]current page.north west) {{\fontsize{12bp}{12bp}\selectfont 部珍宝。”}};
\node[anchor=base west,inner sep=0pt,outer sep=0pt] at ([xshift=72bp,yshift=-644bp]current page.north west) {{\fontsize{12bp}{12bp}\selectfont 【7.17.7】 项羽怒气冲天他说：“明天早晨饱餐士卒，将击溃沛公的军队！”}};
\node[anchor=base west,inner sep=0pt,outer sep=0pt] at ([xshift=72bp,yshift=-660bp]current page.north west) {{\fontsize{12bp}{12bp}\selectfont 【7.17.8】 这时，项羽有兵四十万，驻扎在新丰鸿门，沛公有兵十万，驻扎在霸上。}};
\node[anchor=base west,inner sep=0pt,outer sep=0pt] at ([xshift=72bp,yshift=-676bp]current page.north west) {{\fontsize{12bp}{12bp}\selectfont 【7.17.9】 范增劝告项羽说：“沛公在山东时，贪财好货，喜爱美女。现在进了关，不收}};
\node[anchor=base west,inner sep=0pt,outer sep=0pt] at ([xshift=72bp,yshift=-692bp]current page.north west) {{\fontsize{12bp}{12bp}\selectfont 财物，不亲近妇女，由此看来，他的志向不小。我叫人观望他上空的云气，都呈龙虎形状}};
\node[anchor=base west,inner sep=0pt,outer sep=0pt] at ([xshift=72bp,yshift=-708bp]current page.north west) {{\fontsize{12bp}{12bp}\selectfont ，五颜六色，这是天子之气。赶快进击，不要失掉机会。”}};
\node[anchor=base west,inner sep=0pt,outer sep=0pt] at ([xshift=72bp,yshift=-724bp]current page.north west) {{\fontsize{12bp}{12bp}\selectfont 【7.18.1】 楚国左尹项伯这个人，是项羽的叔父，一向和留侯张良相友好。}};
\end{tikzpicture}
\null
\newpage
\thispagestyle{empty}
\begin{tikzpicture}[remember picture,overlay]
\DrawQuestionTitle
\node[anchor=base west,inner sep=0pt,outer sep=0pt] at ([xshift=72bp,yshift=-84bp]current page.north west) {{\fontsize{12bp}{12bp}\selectfont 【7.18.2】 张良这时跟随着沛公，项伯就夜间骑马跑到沛公军营，私下见到张良，讲述了}};
\node[anchor=base west,inner sep=0pt,outer sep=0pt] at ([xshift=72bp,yshift=-100bp]current page.north west) {{\fontsize{12bp}{12bp}\selectfont 事情的经过，打算叫张良和他一起离去。}};
\node[anchor=base west,inner sep=0pt,outer sep=0pt] at ([xshift=72bp,yshift=-116bp]current page.north west) {{\fontsize{12bp}{12bp}\selectfont 【7.18.3】 他说：“不要跟他们一起死掉。”}};
\node[anchor=base west,inner sep=0pt,outer sep=0pt] at ([xshift=72bp,yshift=-132bp]current page.north west) {{\fontsize{12bp}{12bp}\selectfont 【7.18.4】 张良说：“我为韩王护送沛公，现在沛公的事情发生了危急，逃走是不道义的}};
\node[anchor=base west,inner sep=0pt,outer sep=0pt] at ([xshift=72bp,yshift=-148bp]current page.north west) {{\fontsize{12bp}{12bp}\selectfont ，不能不说一声。”}};
\node[anchor=base west,inner sep=0pt,outer sep=0pt] at ([xshift=72bp,yshift=-164bp]current page.north west) {{\fontsize{12bp}{12bp}\selectfont 【7.18.5】 张良就走了进去，把情况全部告诉了沛公。}};
\node[anchor=base west,inner sep=0pt,outer sep=0pt] at ([xshift=72bp,yshift=-180bp]current page.north west) {{\fontsize{12bp}{12bp}\selectfont 【7.18.6】 沛公大吃一惊，说：“怎么办呢？”}};
\end{tikzpicture}
\null
\end{document}

